\subsection{Schaltungsentwurf eines Operationsverstärkers mit zwei Eingangsspannungen}
% Alt: Aufgabe 11
\Aufgabe{
    \label{Aufg11}
    Zwischen der Ausgangsspannung $U_\mathrm{a}$ und den Eingangsspannungen $U_\mathrm{1}$ und $U_\mathrm{2}$ besteht der folgende
    funktionale Zusammenhang:
    $U_\mathrm{a} = -5 \ U_\mathrm{1} - 0,8 \ U_\mathrm{2}$.
    Geben Sie eine mit Operationsverstärker(n) und Widerständen aufgebaute Schaltung zur Realisierung
    dieses funktionalen Zusammenhangs an. Die Operationsverstärker dürfen als ideale Operationsverstärker
    angenommen werden. Geben Sie Werte der verwendeten Widerstände an!

}

\Loesung{
    Der gegebene funktionale Zusammenhang zwischen der Ausgangsspannung $U_\mathrm{a}$ und den Eingangsspannungen $U_1$ und $U_2$ lautet:
    $
        U_\mathrm{a} = -5 U_1 - 0,8 U_2
    $

    \textbf{1. Schaltungsbeschreibung}

    Zur Realisierung dieses Zusammenhangs nutzen wir eine Summierschaltung mit einem invertierenden Operationsverstärker. Die Ausgangsspannung des invertierenden Summierers lautet:
    $
        U_\mathrm{a} = -\left( \frac{R_\mathrm{F}}{R_1} U_1 + \frac{R_\mathrm{F}}{R_2} U_2 \right)
    $
    Vergleich mit der Vorgabe:
    $
        - \left( \frac{R_\mathrm{F}}{R_1} \right) = -5 \quad \Rightarrow \quad \frac{R_\mathrm{F}}{R_1} = 5
    $
    $
        - \left( \frac{R_\mathrm{F}}{R_2} \right) = -0,8 \quad \Rightarrow \quad \frac{R_\mathrm{F}}{R_2} = 0,8
    $

    \begin{figure}[H]
        \centering
        \input{Tikz/src/LsgOPV2Signale.tex}
        \alttext{Die Abbildung zeigt eine Schaltung mit einem Operationsverstärker, der als invertierender Verstärker mit Rückkopplung ausgeführt ist. Der Operationsverstärker besitzt einen nichtinvertierenden Eingang (+), einen invertierenden Eingang (-) sowie einen Ausgang.
            Der nichtinvertierende Eingang ist direkt mit Masse verbunden. Damit liegt dieser Eingang auf Bezugspotential.
            Der invertierende Eingang ist mit einem Knoten verbunden, an dem mehrere Widerstände zusammengeführt werden. Von diesem Knoten führt der Rückkopplungswiderstand \(R_\mathrm{F}\) zum Ausgang des Operationsverstärkers. Dadurch entsteht eine negative Rückkopplung.
            Zusätzlich sind zwei Eingangssignale angeschlossen:
            Das Signal \(U_{\mathrm{e1}}\) wird über den Widerstand \(R_1\) dem invertierenden Eingang zugeführt.
            Das Signal \(U_{\mathrm{e2}}\) wird über den Widerstand \(R_2\) ebenfalls dem invertierenden Eingang zugeführt.
            Beide Eingangsspannungen werden jeweils gegen Masse gemessen.

            Der Ausgang des Operationsverstärkers befindet sich auf der rechten Seite und ist als Spannung \(U_\mathrm{a}\) gegen Masse eingezeichnet.}
        \label{fig:LsgOPV2Signale}
    \end{figure}

    \textbf{2. Wahl der Widerstandswerte}

    Wir setzen $R_\mathrm{F} = 10\,\mathrm{k}\Omega$ (typischer Wert im Bereich 1–100 k$\Omega$):
    \begin{itemize}
        \item Berechnung von $R_1$:
              $
                  \frac{R_\mathrm{F}}{R_1} = 5 \quad \Rightarrow \quad R_1 = \frac{R_\mathrm{F}}{5} = \frac{10\,\mathrm{k}\Omega}{5} = 2\,\mathrm{k}\Omega
              $
        \item Berechnung von $R_2$:
              $
                  \frac{R_\mathrm{F}}{R_2} = 0,8 \quad \Rightarrow \quad R_2 = \frac{R_\mathrm{F}}{0,8} = \frac{10\,\mathrm{k}\Omega}{0,8} = 12,5\,\mathrm{k}\Omega
              $
    \end{itemize}

    \textbf{3. Realisierung der Schaltung}

    Die verwendeten Widerstände:
    \begin{align*}
        R_\mathrm{F} & = 10\,\mathrm{k}\Omega   \\
        R_1          & = 2\,\mathrm{k}\Omega    \\
        R_2          & = 12,5\,\mathrm{k}\Omega
    \end{align*}


    Siehe: Abschnitt \ref{fig: Verschaltung Verstaerker} und Tabelle \ref{tab:Grundschaltungen}
}